% GENERATED FILE — DO NOT EDIT; authoritative prose is Markdown.
% source: companion/orthability-divine-attributes-and-speech-athari.md
% source-sha256: 64cd7d4b07bd905ee97b16dbffaa046ea51489521e6f06ff47e768b30fe77b53
% source-qualification: research-stage-draft
% source-qualification: not-peer-reviewed
% source-qualification: not-externally-peer-reviewed
% source-qualification: not-empirically-validated
% source-qualification: candidate-terminology-not-adopted
% source-qualification: philosophical-conclusions-conditional-on-stated-premises
% source-qualification: creed-internal-material-school-labeled
% source-qualification: engineering-evidence-does-not-support-metaphysical-claims
\documentclass[10pt,letterpaper,twocolumn]{article}
\usepackage{amsmath}
\usepackage{amssymb}
\usepackage{booktabs}
\usepackage{geometry}
\usepackage{hyperref}
\usepackage{microtype}
\usepackage{natbib}
\usepackage{xcolor}
\geometry{letterpaper,margin=0.75in}
\hypersetup{colorlinks=true,linkcolor=blue,urlcolor=blue,citecolor=blue}
\setlength{\emergencystretch}{3em}
\title{Orthability, the Divine Attributes, and the Uncreated Speech of Allah: An Atharī Creed-Internal Development}
\author{}
\date{}
\begin{document}
\twocolumn[
\maketitle
\textbf{Status: complete creed-internal DRAFT (R6, 2026-07-20; fresh-session repository review completed; not external human peer review; not empirically validated). EXPLICITLY ATHARĪ / school-internal — this paper develops one route within Sunni theology (the Atharī/Ḥanbalī, and specifically Taymiyyan, articulation) and is NOT presented as a criterion-free conclusion of philosophy, nor as the position of every Islamic school.} The Atharī/Taymiyyan route is the declared operative noetic frame; its label identifies premises and custody but never substitutes for warrant. Not peer reviewed. Claim types labeled \emph{analytic}, \emph{transcendental}, \emph{explanatory}, \emph{metaphysical}, \emph{revelational}, or \emph{school-internal}. Philosophical prerequisites are imported from the cross-framework dialectical companion (\href{\detokenize{orthability-and-the-ground-of-intelligibility.md}}{\texttt{orthability-and-the-ground-of-intelligibility.md}}), whose §8 boundary — reason reaches at most a \emph{capacity} for disclosure — governs everything here: every claim beyond it rests on revelation and transmitted creed, and is labeled so. \textbf{Sourcing (R4, Decision 0013):} there is no longer a blanket source-status claim for this paper. \textbf{Each load-bearing claim carries its own machine-readable status} in \href{\detokenize{../references/source-status.yaml}}{\texttt{references/source-status.yaml}} (rows \texttt{ATH-1}…\texttt{ATH-7}), validated in CI. Per-claim statuses are \textbf{not restated here}: the registry rows are authoritative and this front matter defers to them (R5 correction — the previous hand-maintained summary drifted from the registry on ATH-3 and is withdrawn). Two facts worth surfacing at the registry's own recorded strength: the \emph{Majmūʿ al-Fatāwā} vol-12 words/voice formula is carried at \texttt{COMPILATION\_MEDIATED} (row \texttt{ATH-3} — read via the supplied compilation; direct-edition access is the recorded promotion trigger; the adopted doctrine is unchanged and is carried by the primary-text rows), and the courtroom formula is retained \textbf{only} as later-school doctrinal usage, not as history (row \texttt{ATH-6}). Edition-level paginations for six classical loci remain an ordinary research residual (RR-1) with a recorded trigger — not an owner burden. The earlier blanket disclaimer about unopened printed editions is withdrawn: it was inaccurate at the work level (the works were verified) and too coarse at the claim level (only some paginations are edition-dependent). The operational manuscript proves nothing here (absolute firewall), and nothing here is evidence for the operational framework.

\vspace{1em}
]

\section{1. Eight things that must not be collapsed}
The whole discipline of this paper is to keep distinct, and correctly relate, eight items:


\begin{enumerate}
\item \textbf{Allah's eternal Knowledge of all possibilities} — including every determinate linguistic possibility (§4);

\item \textbf{Allah's perfection and attribute of Speech} — He has always been One who speaks, never acquiring the capacity (§3);

\item \textbf{Particular acts of divine Speech occurring when He wills} — His speaking to Mūsā, to Ādam, the sending down of the Qurʾān (§3);

\item \textbf{Allah's uncreated Speech} — His Speech is His attribute, and His attributes are not created (§3);

\item \textbf{The Qurʾān as Allah's uncreated Arabic Speech} — a particular Speech He spoke, in Arabic, uncreated (§3, §5);

\item \textbf{Created human recitation acts} — voices, bodily movements, breath, memorization, and the media: ink, pages, screens (§5);

\item \textbf{Created historical human Arabic} — speakers, conventions, dialects, scripts, grammarians' descriptions (§4);

\item \textbf{The eternally known possibility of those created linguistic instantiations} — item 7 as encompassed, before its existence, by item 1 (§4).

\end{enumerate}
Most classical errors on this topic, and most modern restatements of them, are collapses of two or more of these — the Jahmī/Muʿtazilī collapse of 4–5 into 6–7 ("the Qurʾān is created"), the collapse of 3 into 1 (speech reduced to knowledge), or the reification of 8 into a second eternal entity beside Allah (§4). (\emph{School-internal taxonomy.})


\section{2. From the philosophical companion to the creed: what carries over and what does not}
The cross-framework companion argues to an underived locus only within its stated transcendental scope. Intellect, necessity, Will, Power, Wisdom, and capacity for disclosure remain distinct conditional joints with explicit impersonal, modal, brute-contingency, fittingness, and capacity-sceptical exits; the common-premise fittingness-to-Wisdom bridge is held. Within the Atharī creed, the ground is \textbf{identified} as Allah, and that identification is \emph{revelational}, not a promotion of the cross-framework lane: reason's role is preparatory, while revelation tells us who has spoken and what He said. The creed-internal summary of the philosophical residue, in the owner-approved anti-Platonist form:


\subsection{2.1 Seven Speech bearers}
Predication follows bearer identity. The structured map therefore separates exactly: (1) created human convention; (2) created creaturely speaking, recitation, and writing; (3) created voice and breath; (4) created ink, page, screen, and media; (5) Allah's act of speaking; (6) revealed Arabic wording as Allah's Speech; and (7) created creaturely hearing and reception. Additional recipient-side or historical facts attach to their own bearer. They do not license the unsafe unqualified formula “created Arabic wording,” and the mere capacity for disclosure is not an actual Speech event.


\begin{quote}
Every created act of orthing presupposes an uncreated orthability of being: for any created agent, process, system, or language to resolve, distinguish, signify, classify, or correctly handle anything at all, reality must already be antecedently determinable, differentiable, modally ordered, and fit for lawful resolution. This orthability cannot itself be wholly created, because the intelligible coming-to-be of orthability would already presuppose orthability. In the Atharī framing, that ground is \textbf{not an autonomous Platonic order} but is grounded in Allah's uncreated Knowledge and Speech; created languages and practices instantiate, within creation, possibilities eternally known by Allah. (\emph{Transcendental core + school-internal identification, labeled as such.})

\end{quote}
Sense note (Decision 0010): "orthability" throughout this paper is used in sense \textbf{orthability-O} — the objective conditions of evaluability — never the practice-relative predicate $\operatorname{Orthable}(m; A)$ (sense L) or a created practice's machinery (sense R); the philosophical companion's \href{\detokenize{ARGUMENT-MAP-ORTHABILITY.md}}{\texttt{ARGUMENT-MAP-ORTHABILITY.md}} fixes the senses and the L→O bridge, and nothing in this paper's creed-internal identification changes that typing.


\section{3. The attribute of Speech: the Atharī articulation}

\subsection{3.1 That Allah truly speaks, with speech that is heard}
The creed's basis is textual. "And who is truer in statement than Allah?" (4:87); "and whose words can be truer than those of Allah?" (4:122) — He speaks, and His words are true. "And [remember] when Allah will say: O ʿĪsā son of Maryam…" (5:116) — His speech is heard. "O Mūsā — verily I am your Lord" (20:11–12; the call is 20:11, the declaration 20:12) — words composed of letters. "And We called him from the right side of the Mount, and made him draw near for a talk" (19:52) — calling and conversing occur with sound. And the rule governing all of it: "There is nothing like Him, and He is the All-Hearer, the All-Seer" (42:11) — His speech, letters, and sound do not resemble creatures'. This is the articulation found in Ibn ʿUthaymīn's commentary on \emph{Lumʿat al-Iʿtiqād} (p. 73) \textbf{[via compilation]}: Allah speaks in a real sense, when and as and with whatever He wills, with letters and sound befitting Him. (\emph{Revelational + school-internal.})


\subsection{3.2 Perfection, genus, and particulars: the precise Taymiyyan distinctions}
Four distinctions, each doing real work:


\begin{itemize}
\item \textbf{The perfection is eternal.} Allah has always been capable of speech and has never lacked the perfection of speaking; being capable of speech means He speaks \emph{at will}. (\emph{School-internal; the denial would attribute a deficiency later remedied — excluded.})

\item \textbf{The genus is eternal; the particulars occur when He wills.} In the Taymiyyan formulation, the \emph{kind} — that He speaks — is eternal, while \emph{particular} speeches occur successively by His Will: His speaking to Ādam is a different instance from His speaking to Mūsā; the Qurʾān is a particular Speech He spoke to Jibrīl, spoken at a point after not having been spoken. This is Ibn Taymiyyah's known position on divine speech (developed across \emph{Majmūʿ al-Fatāwā} vol. 12 and elsewhere; academic treatment in Hoover 2004 — whose 2004 article, verified this pass, establishes the eternal-genus/temporal-particulars doctrine for divine \emph{acts and creation}; its extension to divine speech specifically is treated in Hoover's related work, a scope note recorded in the sourcing ledger), and it is what the owner-adopted formulation states: genus eternal → particulars occur → particulars are not beginningless → there is order in the letters and sounds, and internal order excludes a beginningless single utterance. (\emph{School-internal; Taymiyyan/Ḥanbalī articulation specifically — see §7 for the contrasting Sunni account.})

\item \textbf{"Not beginningless" does not mean "created."} A particular occurrence is not thereby a creature: the Qurʾān is a particular Speech, spoken when He willed, \textbf{and it is uncreated}, because speech is an attribute subsisting in the Speaker, and Allah's attributes are not created. The pair ⟨particular occurrence, created thing⟩ must never be identified. (\emph{School-internal; the load-bearing anti-collapse of items 3 and 4 against 6–7.})

\item \textbf{"Eternal" requires disambiguation before use.} The formulae "the Qurʾān is eternal" and "the Qurʾān is not eternal" are both forbidden \emph{unqualified}. Qualified: the perfection and genus of Speech are eternal; the Qurʾān as a particular Speech was spoken when Allah willed (in that sense not a beginningless individual utterance); and it is uncreated. The later Ḥanbalī formula "the word of Allah, ancient and uncreated" (reported of Imām Aḥmad in Ibn Abī Yaʿlā's \emph{Ṭabaqāt al-Ḥanābilah}) uses "ancient" (qadīm) in the sense the school reads as: His Speech as attribute, not a claim that the recited sequence is beginningless. \textbf{Source-critical note (R3):} the \emph{Ṭabaqāt} exists and is correctly attributed (existence verified; page-level locus edition-dependent), but its courtroom narratives are treated in modern scholarship as hagiographically shaped (Cooperson's and Melchert's studies of the Aḥmad biography tradition), and one scholarly line (Madelung's) holds that the "uncreated \textbf{and eternal}" formula hardened only after the miḥna. The formula is therefore cited here \textbf{only} as evidence of how the later school disambiguated "qadīm" — a doctrinal-usage fact — never as verified courtroom history, and no argument in this paper depends on the narrative. (\emph{School-internal doctrinal usage, with the source-critical status stated.})

\end{itemize}

\subsection{3.3 The classical proofs that the Qurʾān is not created}
The tradition's argument set, as compiled from the primary loci \textbf{[via compilation]}, includes: the \textbf{creation/command distinction} — "Surely His is the creation and the command" (7:54), with creation occurring \emph{through} the command "Be!" (36:82), so that if the command were created it would require a prior command, ad infinitum (argued by Sufyān ibn ʿUyaynah and Imām Aḥmad; al-Ājurrī, \emph{ash-Sharīʿah} p. 80; Ḥanbal, \emph{al-Miḥnah} 53–54); the \textbf{knowledge argument} — "The Most Merciful taught the Qurʾān; He created man" (55:1–3), the Qurʾān being from His knowledge (2:120; 3:61), and His knowledge is not created (Aḥmad against al-Qazzāz, \emph{al-Miḥnah} 45); the \textbf{inexhaustibility argument} (18:109; 31:27); the \textbf{names argument} — His names occur in the Qurʾān and are glorified, invoked, and sworn by (87:1; 7:180; 76:25), which would be shirk if created (ath-Thawrī; ash-Shāfiʿī on expiation, \emph{Ādāb ash-Shāfiʿī} 193); the \textbf{attribution argument} — He attributes the revealed Book to Himself as His word (32:2; 6:114; 16:102; 9:6, the last by scholarly consensus about the Qurʾān); the \textbf{seeking-refuge ḥadīth} — "I seek refuge in the perfect words of Allah…" (Muslim 2708), refuge being sought only in Allah and His attributes (Nuʿaym ibn Ḥammād; al-Bukhārī, \emph{Khalq Afʿāl al-ʿIbād} p. 143); the \textbf{superiority ḥadīth} (ad-Dārimī, \emph{ar-Radd ʿalā al-Jahmiyyah} 287, 340); and the \textbf{two rational proofs} — an attribute does not subsist by itself but in the one it describes, so speech attributed to Allah is His attribute and uncreated, whereas the same attribute-kind attributed to a creature is created like the creature (the ninth proof), and the dilemma against created-speech-in-Him or created-speech-outside-Him, the latter yielding the \emph{bush absurdity}: if the words heard by Mūsā were the bush's created words, "Verily I am Allah, Lord of the worlds" (28:30) would stand to Pharaoh's "I am your lord most high" (79:24) with no principled difference — a consequence the tradition treats as self-refuting for the created-speech thesis (the eighth proof; Aḥmad in \emph{al-Miḥnah} 52). (\emph{Revelational + school-internal argumentation, reported with its own loci; this paper adds no claim that these arguments compel the philosophical reader — they are the creed's internal case.})


\section{4. Arabic and the "uncreated grammar" metaphor}

\subsection{4.1 The three-level account (with the fourth level made explicit)}
Before there were Arabic-speaking humans, scripts, grammarians, mouths pronouncing phonemes, or manuscripts in the dunyā, Allah already knew: every possible Arabic sound and form; every root, pattern, inflection, syntactic relation, and semantic distinction; every possible correct and incorrect construction; every context in which an expression would bear one meaning rather than another; and the exact ordered words of the Qurʾān He would speak. Hence the levels:


\begin{enumerate}
\item \textbf{Uncreated divine Knowledge.} Allah's knowledge of Arabic's complete possibility-space is not acquired from created Arabic and does not wait for Arabic to emerge historically: \emph{Arabic's complete orthability is encompassed by what Allah eternally knows.} (\emph{School-internal, resting on the doctrine of His eternal, uncreated Knowledge.})

\item \textbf{The eternal attribute and perfection of Speech} (§3.2's first distinction) — inserted explicitly so that Knowledge and Speech are not collapsed (§6).

\item \textbf{Allah's particular uncreated acts of speaking} — including the Arabic Qurʾān, spoken when He willed, in real ordered words and letters, with a sound befitting Him (§3.2).

\item \textbf{Created Arabic instantiation.} Arabic as historically instantiated \emph{among creatures} — vocal organs, conventional usages, dialect development, writing systems, grammatical descriptions — is created. Creatures do not create the underlying possibility of intelligible linguistic order out of absolute unintelligibility; they instantiate, discover, conventionalize, and modify particular created forms within a possibility-space already known to Allah.

\end{enumerate}
\textbf{Two forbidden sentences, and their corrections.} It is \textbf{wrong to say categorically "Arabic speech is created"}: creaturely Arabic utterances, voices, and media are created; Allah's Arabic Speech — the Qurʾān — is uncreated. The category "Arabic speech" straddles items 5 and 7 of §1 and must be divided before predication. And it is \textbf{wrong to say "letters and sounds existed as created objects before letters and sounds were created"}: nothing of the sort is claimed at level 1 — what "precedes" created Arabic is not a shadow-Arabic of created letters but Allah's Knowledge and His Speech; He never lacked the perfection of Speech, He speaks when He wills, particular speeches occur in an ordered manner, His Speech including the Qurʾān is uncreated, and creaturely phonation and inscription are created. (\emph{School-internal precision rules.})


\subsection{4.2 "Uncreated grammar," defined and defused}
The phrase "uncreated grammar" may be used \textbf{only} as a defined philosophical metaphor: shorthand for \emph{the divine-side ground of expressibility, determinacy, lawful relation, and fitting signification} — that is, for levels 1–2 as they bear on language. It must not be reified as: a second eternal entity; a Platonic Arabic subsisting alongside Allah; or a created human grammar book existing before creation. The possibility-space of language is not a thing beside Allah; talk of it is talk of what Allah eternally knows and of what His Speech is adequate to. Anywhere the metaphor cannot be paraphrased away in that manner, it is being misused. (\emph{School-internal anti-reification rule, matching the anti-Platonism of the philosophical companion.})


\subsection{4.3 The Qurʾān within the three levels}
The Qurʾān is not a human Arabic construction that Allah selected after Arabic independently arose. Its exact wording is Allah's Speech. The created historical existence of Arabic speakers is the condition for \emph{their reception and recitation} of the Arabic Qurʾān — not the ontological source of Allah's ability to speak it: Allah's Speech and Knowledge (uncreated) → the Arabic Qurʾān revealed to creation (heard and read through created media). (\emph{School-internal.})


\section{5. The reciter's act and the carried Word}
When a human recites: the human is created; tongue, mouth, movement, and breath are created; the reciter's voice is created; the act of reciting is created. \textbf{But what is recited — the words — are the words of the Creator.} "The words are the words of the Creator, and the voice is the voice of the reciter" (Ibn Taymiyyah, \emph{Majmūʿ al-Fatāwā} 12/98; similarly 12/53; the Permanent Committee, \emph{Fatāwā al-Lajnah ad-Dāʾimah} 3/21) \textbf{[via compilation]}. Al-Bukhārī devoted \emph{Khalq Afʿāl al-ʿIbād} to exactly this division: the servants' movements, voices, acquisitions, and writing are created; the Qurʾān that is recited, written in the muṣḥaf, and preserved in breasts (29:49; 80:13–14; 98:2–3) is the Word of Allah, uncreated (2/70 and passim; Isḥāq ibn Ibrāhīm: "As for the vessels, who would doubt that they are created?") \textbf{[via compilation]}. Hence it is misleading to say simply "the written Qurʾān is created" or "the heard Qurʾān is created": the ink, page, and voice are created \emph{vessels}; what they carry is the uncreated Word — and the historical Ḥanbalī caution about \emph{lafẓ} formulations (refusing both "my utterance of the Qurʾān is created" and its denial, as ambiguous between act and content) is precisely a refusal to let a vessel-sentence blur the carried-Word distinction. (\emph{School-internal.})


\section{6. Theological precision: the attributes and their distinctions}

\begin{itemize}
\item \textbf{Knowledge} encompasses all possibilities and actualities eternally; it does not come to be. — \textbf{Will} selects among possibilities for creation and for His acts. — \textbf{Power} actualizes what He wills. — \textbf{Wisdom} is the fitting ordering of what He decrees and commands. — \textbf{Speech} is His attribute of address and disclosure: informative and imperative, spoken when He wills.

\item \textbf{Why Speech must not collapse into Knowledge:} the Muʿtazilī reduction ("the Qurʾān expresses His knowledge, and calling it His speech is metaphor") erases item 3 of §1 — the Qurʾānic texts describe address, calling, and audible speaking (5:116; 19:52; 4:164), not mere content-possession. The knowledge \emph{argument} of §3.3 relates the Qurʾān to His knowledge evidentially (it is \emph{from} His knowledge); it does not identify the attributes. Symmetrically, Speech must not collapse into Will-plus-created-sound (the Karrāmī/Jahmī picture of created speech in a substrate — refuted by the bush absurdity, §3.3).

\item \textbf{Why possibility-space must not become an entity:} §4.2's rule, restated at attribute level — the ground of intelligibility is Allah with His attributes, not an abstract order beside Him and not a fifth thing between Creator and creation. (\emph{School-internal.})

\item \textbf{The paper's route-specificity, stated once more:} everything in §§3–6 is the Atharī (and specifically Taymiyyan-inflected) articulation. (\emph{Label.})

\end{itemize}

\section{7. The contrasting Sunni account, described neutrally}
The Ashʿarī and Māturīdī schools affirm — against the Muʿtazilah, and together with the Atharīs — that Allah speaks and that His speech is uncreated. Their articulation differs: Allah's Speech is an eternal attribute subsisting in His essence as \textbf{kalām nafsī} (internal speech/meaning), without letters, sounds, or succession; the recited Arabic Qurʾān is, on the Ashʿarī account, a created \emph{expression} (\textbf{ʿibārah}, al-Ashʿarī's term; the parallel term \textbf{ḥikāyah}, "recounting," belongs to the earlier Kullābī articulation) of that eternal Speech — so that "the Qurʾān is the uncreated Speech of Allah" is affirmed of the kalām nafsī while the articulated wording's status is analyzed differently than in §3. \textbf{The two schools are not identical}, and one classical difference matters here: on the hearing of the divine Speech, Ashʿarī authors (in the Bāqillānī/Juwaynī line, and al-Ashʿarī's own reported position) allow that Mūsā heard the eternal Speech itself, while the Māturīdī position is that the kalām nafsī is not itself heard — what is heard is a created sound by which Allah made him understand it (the difference-lists of the Ibn Kamāl Pāshā genre record this among the classical divergences). Tradition-internal anchors verified this pass: for the Ashʿarī articulation, al-Bājūrī's \emph{Tuḥfat al-Murīd} and al-Bāqillānī's reported formulations; for the Māturīdī creed text, al-Nasafī's \emph{ʿAqāʾid} with al-Taftāzānī's \emph{Sharḥ al-ʿAqāʾid al-Nasafiyyah} — noting that al-Taftāzānī is an Ashʿarī-leaning commentator on a Māturīdī matn. On the Ashʿarī/Māturīdī account the Atharī "letters and sound, when He wills" articulation is rejected as attributing succession (ḥulūl al-ḥawādith) to the divine essence; conversely, Atharī authors object that kalām nafsī makes the revealed Arabic not literally His Speech. \textbf{This paper does not adjudicate that intra-Sunni dispute and does not represent either side as the uncontested meaning of \emph{Ahl as-Sunnah wa-l-Jamāʿah}; it develops the Atharī route, and this paragraph exists so the reader cannot mistake the route for the whole map.} (\emph{Neutral description; sources SECONDARY-VERIFIED at the tradition-internal level in R3 — see the companion sourcing ledger; page-level edition checks remain a recorded residual.})


\section{8. Reason and revelation: the final accounting}

\begin{center}
\begin{tabular}{@{}p{0.470\linewidth}p{0.470\linewidth}@{}}
\toprule
\textbf{Claim} & \textbf{Basis} \\
\midrule
An underived ground of intelligibility, not intrinsically incapable of disclosure & Philosophical (transcendental + conditional; cross-framework dialectical companion, with its stated exits) \\
That the ground is Allah, who has truly spoken & Revelational \\
That His Speech is uncreated; letters and sound befitting Him; heard by whom He wills & Revelational + school-internal (Atharī articulation, §3) \\
Genus eternal; particular speeches when He wills; particular ≠ created & School-internal (Taymiyyan precision, §3.2) \\
The Qurʾān is His uncreated Arabic Speech; the reciter's voice and the media are created & Revelational + school-internal (§5) \\
Arabic's possibility-space eternally known; created Arabic instantiated within it; no reified "uncreated grammar" & School-internal synthesis (§4), consistent with — not proven by — the philosophical companion \\
No engineering result supports any row of this table & Firewall (absolute) \\
\bottomrule
\end{tabular}
\end{center}
Reason, at its strongest here, removes alleged impossibilities and exhibits presuppositions. It does not recite. What the Qurʾān is — Allah's Word, uncreated, "and if anyone of the polytheists seeks your protection, grant it, that he may hear the Word of Allah" (9:6) — is known from revelation, held as creed, and articulated by this school in the distinctions this paper has tried to keep exact. (\emph{Closing; school-internal.})


\bigskip\hrule\bigskip

\section{References}
Sources actually cited in this paper. Bibliographic records live in \href{\detokenize{../references/orthemology.bib}}{\texttt{references/orthemology.bib}}. \textbf{Per-claim evidence-access statuses are NOT restated here} — they live in \href{\detokenize{../references/source-status.yaml}}{\texttt{references/source-status.yaml}} rows \texttt{ATH-1}…\texttt{ATH-7} and the sourcing surfaces (start at \href{\detokenize{../docs/sourcing/CURRENT-SOURCING-LEDGER.md}}{\texttt{docs/sourcing/CURRENT-SOURCING-LEDGER.md}}); a work's presence in this list asserts bibliographic identity only, never a wording-access claim. Classical works are cited with the school-internal loci given inline; several page loci are edition-dependent (research residual RR-1).


\begin{itemize}
\item \emph{al-Qurʾān al-Karīm} (loci per \href{\detokenize{../references/quran-loci.yaml}}{\texttt{references/quran-loci.yaml}}, primary-verified).

\item Muslim ibn al-Ḥajjāj. \emph{Ṣaḥīḥ Muslim} (ḥadīth 2708).

\item al-Bukhārī, Muḥammad ibn Ismāʿīl. \emph{Khalq Afʿāl al-ʿIbād}.

\item Ibn Taymiyyah, Aḥmad. \emph{Majmūʿ al-Fatāwā} (vol. 12; the words/voice formula is carried at registry row \texttt{ATH-3}'s recorded strength).

\item Ibn ʿUthaymīn, Muḥammad ibn Ṣāliḥ. \emph{Sharḥ Lumʿat al-Iʿtiqād}.

\item al-Ājurrī, Muḥammad ibn al-Ḥusayn. \emph{ash-Sharīʿah}.

\item Ḥanbal ibn Isḥāq. \emph{Dhikr Miḥnat al-Imām Aḥmad ibn Ḥanbal}.

\item ad-Dārimī, ʿUthmān ibn Saʿīd. \emph{ar-Radd ʿalā al-Jahmiyyah}.

\item Permanent Committee for Scholarly Research and Iftāʾ. \emph{Fatāwā al-Lajnah ad-Dāʾimah} (vol. 3).

\item Ibn Abī Ḥātim ar-Rāzī. \emph{Ādāb ash-Shāfiʿī wa-Manāqibuh}.

\item Hoover, Jon (2004). "Perpetual Creativity in the Perfection of God: Ibn Taymiyya's Hadith Commentary on God's Creation of this World." \emph{Journal of Islamic Studies} 15(3).

\item Cooperson, Michael (2000). \emph{Classical Arabic Biography: The Heirs of the Prophets in the Age of al-Maʾmūn}. Cambridge University Press.

\item Melchert, Christopher (2006). \emph{Ahmad ibn Hanbal}. Oneworld.

\end{itemize}

\twocolumn
\bibliographystyle{plainnat}
\bibliography{../../../references/orthemology}
\end{document}
